% SHARD-ID: CA-102-RECORD-FORGETS-THE-POINT
% SHARD-TITLE: The Interval/Dilation Record Forgets the Point
% SHARD-SUMMARY: The interval/dilation record does not determine the pointed input: charge translation in an abelian graded family preserves every interval and the fixed-H refinement while changing the site class, refuting the positive form of Target 97.16.
% SHARD-SUMMARY: Cut commutants are intrinsic, but vacancy projections must be carried; even named points do not make placement maps natural. Pointed records separate the tested counterexamples, while general pointed-record faithfulness remains open.
% SHARD-KEYWORDS: vacancy, pointed record, interval algebra, dyadic refinement, relative commutant, charge translation, half-braiding, naturality, fidelity, P1, G6
\section{The Interval/Dilation Record Forgets the Point}\label{sec:record-forgets-the-point}

\subsection{Status, Sources, and Dependencies}

\textbf{Status: COUNTEREXAMPLE} to the positive form of CA-97 Target 97.16, exactly as stated.
Statements 102.1, 102.3, 102.5--102.7 are \emph{derived}; 102.2 and 102.4 are \emph{negative};
102.8 is \emph{conditional}, an unproved extension with a Law-1 gap; 102.9 is \emph{conjectural}.
NOTE, REVIEW and INTEGRATION mean the files in
\path{explorations/2026-10-25-astra-r3/G6-record-faithfulness/}.
\textbf{Pillar P1:} the vacancy in \(O\cong\one\oplus X\), and hence the indefinite-number interpretation,
must be carried along the tower as declared structure. \textbf{P4 interface:} this sharpens the
\(\mathsf R4\) insufficiency annotation; it supplies no continuum certificate.
\textbf{CA-88 cells (INTEGRATION:36--46):} (any/open; \(\mathsf L0+H\); \(\mathsf R0^u\)), dyadic geometry,
gains a negative information annotation, not an existence or attainment change. The \(\mathsf R4\)
insufficiency annotation extends to the full interval/dilation record: the first missing mark proved here
is the point. Trailing \((\mathrm{any};\mathsf L0;\mathsf R0^u)\), \(\mathsf R0^c\), \(\mathsf R1\)--\(\mathsf R5\)
attainment, full/periodic and \(\mathsf L3\) cells do not move; CA-93's fidelity hypotheses remain required.
All settings below are \emph{any, ordered open geometry}; no chiral/full continuum target is selected.
Dependencies: CA-77 (77.2), CA-88 (grid), CA-90 (90.H, 90.1), CA-97 (97.11, 97.13, 97.16);
CONVENTIONS (a), (b), (r), (am), (aq), (as), (ax). Additional notation is scoped here;
its proposed registration is in the package's \path{SHARD_META.md}.

% Source: explorations/2026-10-25-astra-r3/G6-record-faithfulness/NOTE.md:17-35,43-99,101-200,202-296 -- record and proofs.
% Source: explorations/2026-10-25-astra-r3/G6-record-faithfulness/REVIEW.md:3-19,23-91,99-177,179-224 -- grades, corrections, R-1--R-3, target.
% Source: explorations/2026-10-25-astra-r3/G6-record-faithfulness/INTEGRATION.md:12-46,48-88,90-128 -- lift, cells, amendments, notation, target.
% Source: explorations/2026-10-25-astra-r3/briefs/G6-record-faithfulness.md:83-118 -- positive conjecture and requested tests.
% Source: report/sections/77_pointed_carrier_invariants.tex:47-81 -- point and sufficient marked reconstruction.
% Source: report/sections/90_local_unital_lift.tex:59-77,110-141 -- H and local dyadic shuffle.
% Source: report/sections/97_what_the_tower_remembers.tex:178-218,265-272 -- fixed-H descent, typed extraction, exact target.
% Source: report/sections/88_layered_targets.tex:96-125,149-154 -- cells and fidelity boundary.
% Source: report/sections/93_normal_limit_expectations_and_fidelity.tex:158-177 -- fidelity requires a conformal net and certificates.
% Source: references/text/PenneysUnitaryFusionCategories.md:219-237,277-288,319-348,394-398,470-478 -- splitting, graded categories, tensor functors, coherence.
% Source: references/text/EtingofNikshychOstrikFusionCategories.txt:1459-1487,1700-1706 -- positive dimensions, unitary bridge.
% Source: references/manifest/SOURCES.md:34-35,55-57 -- local source registration.
% Source: CONVENTIONS.md:1450-1470,1524-1543,1564-1575,1644-1665 -- point, tags, crossing, deferred registration.
% Source: explorations/2026-10-25-astra-r3/G6-record-faithfulness/checks/check_translation.py:1-137 -- exact matrix-unit calibration.
% Source: explorations/2026-10-25-astra-r3/G6-record-faithfulness/checks/translation_output.txt:1-11 -- exact calibration output.
% Source: explorations/2026-10-25-astra-r3/G6-record-faithfulness/checks/review_g6_checks.py:1-31,186-268 -- phase-sensitive and point-block checks.
% Source: explorations/2026-10-25-astra-r3/G6-record-faithfulness/checks/review_g6_output.txt:1-13 -- GREEN phase and point-block checks.
% Source: explorations/2026-10-25-astra-r3/G6-record-faithfulness/checks/review_g6_mutations.txt:1-12 -- rejected mutations.
% Source: worklog/023_2026-10-27_status_overview.md:179-188 -- current P1/fidelity diagnosis; scope correction below.
Checks: the exact matrix-unit script and independent NumPy reviewer script reproduce their stored outputs;
the latter checks crossing phases and point-block sizes. Their finite checks do not prove the all-length statements.

\subsection{The General Record and Cut Information}

Fix arbitrary unitary fusion \(\Cc\), pointed \((O,\eta)\), and \(\mathcal D=\langle O\rangle\subset\Cc\).
Let \(H=(t_Y)\) be one fixed pointed unitary relative half-braiding on \(\mathcal D\) (90.H):
\(t_Y:Y\otimes O\to O\otimes Y\) is natural, tensorial and unital, with
\(t_Y(1_Y\otimes\eta)=\eta\otimes1_Y\). Put \(W_L=O^{\otimes L}\), \(W_0=\one\),
\(A_L=\operatorname{End}_{\Cc}(W_L)\), and let \(\kappa_L:A_L\to A_\infty\) be the norm-limit maps
for trailing inclusion \(a\mapsto a\otimes1_O\). Bracketing is left-associated; fixed coherent associators are retained.
The shuffle \(B_L\) crosses old \(j\) with new \(i\) exactly when \(i<j\), and
\(\alpha_L(a)=B_L(a\otimes1_{W_L})B_L^\dagger\), \(\delta\kappa_L=\kappa_{2L}\alpha_L\) (97.11).
The record and its proposed pointed augmentation are
\[
 \mathcal R=(A_\infty,\{A_{[j,k]}\},\delta),\quad
 A_{[j,k]}=\kappa_k(1_{W_{j-1}}\otimes A_{k-j+1}),\qquad
 \mathcal R^p=(\mathcal R,\{p_j\}),\quad p_j=\kappa_j(1_{W_{j-1}}\otimes\eta\eta^\dagger).
\]
An isomorphism preserves exactly the listed marks, every interval label, and \(\delta\).
The \(p_j\) are placed copies of (am)'s \(p_0\), not its occupation projection \(p_1\).
``Natural recovery'' means transport by \emph{every} record isomorphism; recovery up to a coherent choice is weaker.

\textbf{Lemma 102.1 [derived; any/open; \(\mathsf L0+H\); \(\mathsf R0\) finite information]\newline (NOTE:43--99; REVIEW:70--73).}
For any cut \(m|n\), put \(V_m^x=\Hom(x,W_m)\), \(F_{xy}^z=\Hom(z,x\otimes y)\); omit zero summands.
Under orthogonal simple splittings,
\begin{equation}\label{eq:102-cut}
\begin{split}
 A_{m+n}&\cong\bigoplus_z\operatorname{End}\Big(\bigoplus_{x,y}V_m^x\otimes V_n^y\otimes F_{xy}^z\Big),\\
 (A_{[1,m]}\vee A_{[m+1,m+n]})'\cap A_{m+n}
 &\cong\bigoplus_{z,x,y}1_{V_m^x\otimes V_n^y}\otimes\operatorname{End}(F_{xy}^z).
\end{split}
\end{equation}
The record preserves these commutants, their inclusion diagrams, and intersections involving \(\delta^r(A_I)\).
\par\(\langle1\rangle1\). Compose orthonormal isometric splittings of \(W_m,W_n,x\otimes y\)
(Penneys:219--237,319,470--478); the two interval algebras act on the two \(V\) factors.
Their product embedding is faithful: \(d_xd_y>0\), so \(x\otimes y\ne0\)
(ENO:1459--1487,1700--1706, including the reviewer's missing unitary bridge).
\par\(\langle1\rangle2\). Within each \(z\)-block, the two central projections force preservation of \((x,y)\).
Commuting with matrix units kills off-diagonal entries on \(V_m^x\otimes V_n^y\) and equalises diagonal entries;
the remaining operator is \(1\otimes c\). Conversely every such operator commutes.
\par\(\langle1\rangle3\). Isomorphisms preserve multiplication, commutation and intersections; intertwining
\(\delta\) gives the moved-cut assertion. QED. The formula alone uses \(\mathsf L0\); the last assertion needs \(H\).
The displayed simple labels are choices, not intrinsically named charges.

\subsection{An Exact Counterexample Family}

\textbf{Negative 102.2 [negative; any/open; \(\mathsf L0+H\); \(\mathsf R0^u\) information]\newline (NOTE:101--184; REVIEW:74--79,170).}
The universal claim that \(\mathcal R\) determines the pointed input for arbitrary \(\Cc\) is false. The mechanism is
proved for the subclass \(\mathcal D=\Hilb_{\mathrm{fd}}(G,1)\), finite abelian \(G\), with trivial associator.
Write \(O=\bigoplus_{a\in G}\mathbb C^{m(a)}\otimes a\). Choose \(s\in G\) with \(m(e),m(s^{-1})>0\);
both \(\operatorname{supp}m\) and its translate must generate \(G\). Let \(O'=sO\) retain multiplicity spaces
but replace grade \(a\) by \(sa\); choose unit points in \(O_e,O'_e\), and initially take flip \(H\) on both.
Then \(\mathcal R(O)\cong\mathcal R(O')\), even preserving translations and endomorphism concatenation.
If no group automorphism matches the two multiplicity functions, no tensor equivalence carries \(O\) to \(O'\).
\par\(\langle1\rangle1\). The flip is grade-preserving since \(G\) is abelian, natural on graded maps,
unitary and pointed; its tensor identity sends \(y\otimes z\otimes v\) to \(v\otimes y\otimes z\)
in either order. This verifies 90.H in the sourced graded construction (Penneys:277--288,319--321).
\par\(\langle1\rangle2\). The underlying Hilbert unitary \(U:O\to O'\) preserves multiplicity bases and
has degree \(s\); it is \emph{not} a categorical morphism for \(s\ne e\). Put \(U_L=U^{\otimes L}\),
\(\phi_L=\operatorname{Ad}U_L\), \(\phi_0=\mathrm{id}\). A word of total charge \(x\) acquires charge \(s^Lx\).
Thus \(\phi_L\) maps all categorical matrix units bijectively and is an isometric unital \(*\)-isomorphism.
\par\(\langle1\rangle3\). Factorwise conjugation gives simultaneously, for \(b\in A_{k-j+1}\),
\[
 \phi_{L+1}(a\otimes1)=\phi_L(a)\otimes1,\qquad
 \phi_L(1^{\otimes(j-1)}\otimes b\otimes1^{\otimes(L-k)})
 =1^{\otimes(j-1)}\otimes\phi_{k-j+1}(b)\otimes1^{\otimes(L-k)}.
\]
These are onto on every interval; also \(\phi_{L+M}(a\otimes b)=\phi_L(a)\otimes\phi_M(b)\).
\par\(\langle1\rangle4\). The flip shuffle is a tensor permutation, hence
\begin{equation}\label{eq:102-intertwiner}
 U_{2L}B_L=B'_L(U_L\otimes U_L),\qquad
 \phi_{2L}\alpha_L=\alpha'_L\phi_L,\qquad \Phi\delta=\delta'\Phi.
\end{equation}
For the last identity extend the compatible \(\phi_L\) and their inverses to the norm completions and use
\(\Phi\delta\kappa_L=\kappa'_{2L}\phi_{2L}\alpha_L=\delta'\Phi\kappa_L\), then density (97.11).
\par\(\langle1\rangle5\). A tensor equivalence permutes invertible simples by a group automorphism
and preserves multiplicities (Penneys:325--348); the stated mismatch forbids it. QED.
In \eqref{eq:102-cut}, \(F_{xy}^z\) has dimension one exactly when \(z=xy\); the relabelling
\((x,y,z)\mapsto(s^mx,s^ny,s^{m+n}z)\) preserves every cut diagram, including \(\delta\)-moved cuts.
This demonstrates insufficiency for this family, not a reconstruction obstruction for every \(\Cc\).

\textbf{Proposition 102.3 [derived; any/open; \(\mathsf L0+H\); \(\mathsf R0^u\) information]\newline (NOTE:154--177; REVIEW:106--121).}
Reviewer R-1 extends 102.2 to \(t_Y(y_a\otimes v_b)=\chi(a,b)v_b\otimes y_a\), with unitary
bicharacter \(\chi\), provided \(\chi(a,s)=\chi(s,a)=1\) for every \(a\in G\).
\par\(\langle1\rangle1\). Bicharacter identities give 90.H, including \(\chi(a,e)=1\) at the point.
\par\(\langle1\rangle2\). Transparency gives \(\chi(sa,sb)=\chi(a,b)\) at each crossing, proving
\eqref{eq:102-intertwiner} with the same phases. The other steps are unchanged. QED.
Failure of this particular \(\Phi\) for a nontransparent \(s\) excludes no other possible record isomorphism.

\textbf{Negative 102.4 [negative; any/open; \(\mathsf L0+H\); \(\mathsf R0\) information]\newline (NOTE:186--200; REVIEW:47--63,80--81).}
With flip \(H\), let \(R_G=\bigoplus_{a\in G}a\).
\emph{B1:} for every nontrivial finite abelian \(G\) and \(s\ne e\), take
\(O=R_G\oplus\one\), \(O'=R_G\oplus s\). Both supports generate; translation gives 102.2, but
\(\dim\Hom(\one,O)=2\ne1=\dim\Hom(\one,O')\), invariant under every tensor equivalence.
\emph{B2:} take nonunit \(a,b\) in distinct automorphism orbits, \(s=ba^{-1}\),
\(O=R_G\oplus a\), \(O'=R_G\oplus b\). Both are connected; an equivalence would carry the uniquely
doubled simple \(a\) to \(b\). Different orders suffice; this order-based witness requires such elements in \(G\).
Thus the positive form of 97.16 is false, including for connected sites. For \(s\ne e\),
\(\phi_1(\eta\eta^\dagger)\) lies in grade \(s\), whereas the target point lies in grade \(e\).
The record forgets the point and the marks built from it; \(\mathcal D\) and the supplied transport rule are unchanged.

\subsection{Named Points and the Remaining Marks}

\textbf{Proposition 102.5 [derived; any/open; \(\mathsf L0+H\) and (G6.5); \(\mathsf R0\) marks]\newline (NOTE:202--269; REVIEW:82--85).}
For arbitrary \(\Cc\) in the general setting, assume a noncentral unitary \(u\in\operatorname{End}(O)\) with
\[
 u\eta=\eta,\qquad t_Y(1_Y\otimes u)=(u\otimes1_Y)t_Y\quad(Y\in\mathcal D). \tag{G6.5}
\]
Then \(\mathcal R^p\) has an automorphism moving the standard translation, concatenation,
left-insertion corner and hop \(h\), although their relevant range projections stay fixed.
\par\(\langle1\rangle1\). Put \(C_0=C_1=1\), \(C_L=1_O\otimes u^{\otimes(L-1)}\),
\(\gamma_L=\operatorname{Ad}C_L\). Since \(C_{L+1}=C_L\otimes u\) for \(L\ge1\), these maps
respect trailing inclusion and map every interval algebra onto itself; extend to \(\gamma\) by norm density.
\par\(\langle1\rangle2\). Naturality transports \(u\) on old strands, (G6.5) on ancillary strands;
old strand 1 never crosses. Thus \(B_L(C_L\otimes u^{\otimes L})=C_{2L}B_L\).
Conjugate \(a\otimes1\): ancillary factors cancel, so \(\gamma_{2L}\alpha_L=\alpha_L\gamma_L\)
and \(\gamma\delta=\delta\gamma\). Also \(\gamma(p_j)=p_j\) since \(u\eta=\eta\).
\par\(\langle1\rangle3\). Choose \(b\) with \(ubu^\dagger\ne b\). For \(T(b)=1\otimes b\),
\(\gamma_2T(b)=1\otimes ubu^\dagger\ne T\gamma_1(b)\); the same inequality moves concatenation
\((a,b)\mapsto a\otimes b\), and \(\theta(b)=p_0\otimes b\), \(p_0=\eta\eta^\dagger\).
The latter range \(p_0\otimes1\) is fixed, and compression by \(\eta\) proves the inequality persists.
\par\(\langle1\rangle4\). For \(r=\eta\otimes1,l=1\otimes\eta:O\to W_2\), \(h=rl^\dagger\),
\(C_2r=ru\), \(C_2l=l\), hence \(\gamma_2(h)=rul^\dagger\ne h\), with both support projections fixed.
\par\(\langle1\rangle5\). Every input admits such an instance after replacing \(O\) by \(O\oplus\one\)
and extending \(H\) by the unit transport. Set \(u=+1\) on the chosen point, \(-1\) on the new unit copy,
and identity elsewhere. Pointedness and unitarity split off the old unit transport; (G6.5) holds blockwise,
while a matrix unit between the two unit copies witnesses noncentrality. Generation is unchanged. QED.
This proves \emph{non-naturality}, not failure of reconstruction up to a coherent choice.

\textbf{Proposition 102.6 [derived; any/open; \(\mathsf L0+H\); \(\mathsf R0\) information]\newline (NOTE:281; REVIEW:146--158).}\par
For arbitrary inputs, an \(\mathcal R^p\)-isomorphism preserves the matrix-block size containing
\(p_1\cdots p_L\) for each \(L\). Consequently it separates the three tested B1/B2 pairs in the table below
(reviewer R-3, derived for these instances only).
\par\(\langle1\rangle1\). The product is the rank-one projection onto \(\eta^{\otimes L}\)
in the unit block of \(A_L\), of size \(m_L(\one)=\dim\Hom(\one,W_L)\).
\par\(\langle1\rangle2\). A \(*\)-isomorphism carries its minimal central support to the central support
of the image and preserves that full matrix algebra's size. The displayed unequal sizes prohibit every
point-preserving isomorphism. QED. No general sufficiency theorem for \(\mathcal R^p\) follows.

\textbf{Proposition 102.7 [derived; any/open; \(\mathsf L0+H\); \(\mathsf R0\) inventory]\newline (NOTE:271--296; REVIEW:85--88,173--174).}
The record intrinsically retains named stages, inclusions, centres, block sizes, commutants and products of
\emph{already embedded} observables (definition and 102.1). The corrected refinement entry is only
\(\delta:A_{[j,k]}\hookrightarrow A_{[2j-1,2k-1]}\) (90.1(ii)); the reviewed GAP was its upgrade to an identification.
Given \(p_j\), insertion-range projections are products over omitted sites; their domain identifications are
extra natural marks (102.5). The point itself is extra (102.4); concatenation alone cannot repair its loss (102.2).
CA-77 Proposition 77.2's coherent placements and partial isometries are a sufficient richer record;
97.13 additionally recovers supplied \(H\) from typed compressions. No minimal sufficient list is proved.

\textbf{Proposition 102.8 [conditional; any/open; \(\mathsf L0+H\); \(\mathsf R0\) information]\newline (NOTE:40--41; REVIEW:123--136).}
Reviewer R-2 is an \emph{unproved conjectural sketch}, kept conditional here: for arbitrary unitary fusion
\(\mathcal E\) with pointed generating \((O_{\mathcal E},\eta_{\mathcal E},H_{\mathcal E})\), a proposed
\(\mathcal E\boxtimes\Hilb_{\mathrm{fd}}(G)\) realisation would use
\(O=O_{\mathcal E}\boxtimes O_G\) and \(\mathrm{id}\otimes\operatorname{Ad}U_L\).
It requires a sourced construction and compatible factorisation of endomorphisms, intervals and shuffles;
no local source for that product/graded-object construction was supplied. This is not a proved extension
to every \(\Cc\) possessing an invertible object, and is not used in 102.2--102.7.

\subsection{Effect on other shards}
\emph{Amendment pointers (2026-10-27, round 4; INTEGRATION:48--88).} CA-97's status, summary and
97.16 must record the negative 102.4 and the replacement 102.9 (97.16\('\)); after 97.13, name the point as
a necessary carried mark. CA-77's placements remain part of its sufficient record (102.2, 102.5);
the converse to CA-90 90.1(iv) fails for \(\mathcal R\) (102.4). CA-88 gains only the information annotations above.
Worklog 023 \S4.4 must limit G6-B's proved mechanism to the abelian graded family with trivial associator
and compatible \(H\): the same category occurs on both sides, and pointed-record recovery of \(F,H\) remains open.
The pipeline must carry named vacancy projections, or richer coherent placement data determining them.
Exact append-after anchors and dated amendment text are supplied in \path{SHARD_META.md}; these files are not edited here.

\subsection{What This Does Not Claim}
There is no proof that \(\mathcal R\) forgets \(F\) or \(H\) once points are named, and no general charge-translation
classification. G6-D's different-associator and different-transport tests remain open (REVIEW:87--89);
the nonexistent G6-D citation is discarded. Non-natural marks may still admit coherent choices (102.5).
No state, Wilson limit, symmetry generators, conformal limit or \(\mathsf R4\) certificate is constructed.

\subsection{Examples (Instances Only)}
All rows instantiate the statements above with (any/open; \(\mathsf L0+H\); \(\mathsf R0\) information).
Grades are additive, unit \(0\); displayed site lists include multiplicities. Points are unit basis vectors
(the first unit copy when repeated); \(\chi=1\) unless specified. Block sizes and failures are the checked
data in \path{checks/review_g6_output.txt}:1--13, not claims for arbitrary \(\Cc\).
\begin{center}\footnotesize
\begin{tabular}{p{2.1cm}p{12.0cm}}\toprule
Instance & Exact data and verdict \\\midrule
B1, \(\mathbb Z_2\) & \(O=(0,0,1)\), \(O'=(0,1,1)\), \(s=1\): 102.4 [negative]; point-block sizes \((2,1)\) at \(L=1\), 102.6 [derived]. \\
B2, \(\mathbb Z_4\) & \(O=(0,1,1,2,3)\), \(O'=(0,1,2,2,3)\), \(s=1\); doubled orders \(4,2\): 102.4 [negative]; point-block sizes \((6,7)\) at \(L=2\) [derived]. \\
B2, \(\mathbb Z_6\) & \(O=(0,1,1,2,3,4,5)\), \(O'=(0,1,2,2,3,4,5)\), \(s=1\); doubled orders \(6,3\): 102.4 [negative]; point-block sizes \((57,58)\) at \(L=3\) [derived]. \\
R-1, \(\mathbb Z_4\) & \(O=(0,0,1,2,3)\), \(O'=(0,1,2,2,3)\), \(s=2\), \(\chi(a,b)=(-1)^{ab}\): transparent, 102.3 [derived]. Plain \(\Phi\) fails dilation at \(L=2\) for this site and \(\chi=i^{ab}\); also for the B2 site, \(s=1\), \(\chi=(-1)^{ab}\) or \(i^{ab}\). \\
C1, \(\mathbb Z_2\) & \(O=(0,1,0)\), \(\eta=(1,0,0)\), \(u=\operatorname{diag}(1,1,-1)\): 102.5 [derived], \(\gamma\) fixes points and moves \(T,h\). \\\bottomrule
\end{tabular}
\end{center}

\subsection{Single Next Target: Pointed-Record Faithfulness}

\textbf{Target 102.9 [conjectural; any/open; \(\mathsf L0+H\); \(\mathsf R0\) information only]\newline (NOTE:319--323; REVIEW:208--224).}
This is the reviewer's 97.16\('\) (INTEGRATION:117--128). For arbitrary pointed tensor-generating
\((\mathcal D,O,\eta,H)\) in a unitary fusion \(\Cc\), does \(\mathcal R^p\cong\mathcal R'^p\) force
a unitary tensor equivalence \(\mathcal D\simeq\mathcal D'\) transporting \(O,\eta,H\)? A positive route must correct the isomorphism by
a point-preserving automorphism so that coherent placements, \(h\) and typed compressions \(i,j\) are
transported up to unitary gauge; then \(t_O=j^\dagger\alpha_2(h)i\) (97.13) and 77.2 apply.
A negative route requires one simultaneous point-, interval- and dilation-preserving counterexample.
Neither the unpointed intertwiner 102.2 nor the non-naturality result 102.5 settles this question.
